\subsection*{CU-13: Cambiar el nickname}
\addcontentsline{toc}{subsection}{CU-13: Cambiar el nickname}
\label{cu-13}

\begin{fichacu}{CU-13}{Cambiar el nickname}
  \crudFila{Responsable}{Ángel Gabriel Aguilar Hernández}
  \crudFila{Actor principal}{Jugador}
  \crudFila{Actores de sistema}{Servidor de partidas, Base de datos}
  \crudFila{Prototipos}{9c, 9d}
  \crudFila{Objetivo}{Permitir al Jugador cambiar una única vez el nombre con el que lo ven los demás, dejando constancia del cambio para que su historial siga siendo rastreable.}
  \crudFila{Disparador}{El Jugador selecciona ``Cambiar'' junto al \texttt{nickname} en la \texttt{GUI\_\allowbreak{}EditProfile}.}
  \textbf{Precondiciones} & \cuCelda{\begin{cureglas}
      \item \textbf{\mbox{PRE-01}} El Jugador tiene una \textit{Sesion} abierta y su token es válido.
      \item \textbf{\mbox{PRE-02}} El \textit{Usuario} tiene \texttt{tipo\_\allowbreak{}cuenta} en \texttt{REGISTRADA}. Una cuenta de invitado no puede gastar el cambio: su \texttt{nickname} es provisional y se elige al vincular (\mbox{RN-06}).
      \item \textbf{\mbox{PRE-03}} El \textit{Usuario} no ha cambiado su \texttt{nickname} con anterioridad.
      \item \textbf{\mbox{PRE-04}} El Servidor de partidas está en ejecución y tiene conexión disponible con la Base de datos.
    \end{cureglas}} \\ \hline
  \textbf{Flujo normal} & \cuCelda{\begin{cuenum}[start=1]
      \item El SJM despliega la \texttt{GUI\_\allowbreak{}ChangeNickname} con el \texttt{nickname} actual, un campo para el nuevo y un aviso destacado de que el cambio es \textbf{irreversible y único} (\mbox{RN-01}). \textit{(Ver \mbox{FA-01})}
      \item El Jugador captura el \texttt{nickname} nuevo y da click en ``Continuar''.
      \item El SJM valida que mida entre tres y treinta caracteres y que no contenga espacios ni caracteres no admitidos (\mbox{RN-02}). \textit{(Ver \mbox{FA-02})}
      \item El SJM valida que sea distinto del actual. \textit{(Ver \mbox{FA-03})}
      \item El SJM despliega la \texttt{GUI\_\allowbreak{}MessageConfirm} con el texto exacto del \texttt{nickname} nuevo y la advertencia de que no habrá otra oportunidad, con las opciones ``Cambiar'' y ``Cancelar'' (\mbox{RN-07}). \textit{(Ver \mbox{FA-04})}
      \item El Jugador selecciona ``Cambiar''.
      \item El SJM envía el mensaje \texttt{CHANGE\_\allowbreak{}NICKNAME\_\allowbreak{}REQUEST} con el \texttt{nickname} nuevo y el token de sesión. \textit{(Ver \mbox{EX-03}, \mbox{EX-04})}
      \item El Servidor de partidas verifica que el token corresponda a una \textit{Sesion} abierta. \textit{(Ver \mbox{FA-05})}
    \end{cuenum}} \\
   & \cuCelda{\begin{cuenum}[start=9]
      \item El Servidor de partidas verifica que el \textit{Usuario} no tenga ningún \textit{HistorialNickname}, es decir, que no haya gastado su cambio (\mbox{RN-01}). \textit{(Ver \mbox{FA-06})}
      \item El Servidor de partidas vuelve a aplicar las validaciones de los pasos 3 y 4 (\mbox{RN-08}). \textit{(Ver \mbox{FA-07})}
      \item El Servidor de partidas verifica que el \texttt{nickname} no contenga términos del catálogo de \textit{PalabraProhibida} con ámbito de nickname. \textit{(Ver \mbox{FA-08})}
      \item El Servidor de partidas verifica que el \texttt{nickname} no esté en uso por otro \textit{Usuario}, comparando sin distinguir mayúsculas ni acentos (\mbox{RN-02}). \textit{(Ver \mbox{FA-09})}
      \item El Servidor de partidas inicia una transacción sobre la Base de datos.
      \item El Servidor de partidas crea el registro de \textit{HistorialNickname} con el \texttt{nickname} anterior, el nuevo y la fecha del cambio (\mbox{RN-03}). \textit{(Ver \mbox{EX-01})}
      \item El Servidor de partidas actualiza el \texttt{nickname} del \textit{Usuario}.
      \item El Servidor de partidas confirma la transacción. \textit{(Ver \mbox{EX-02})}
    \end{cuenum}} \\
   & \cuCelda{\begin{cuenum}[start=17]
      \item El Servidor de partidas responde con \texttt{CHANGE\_\allowbreak{}NICKNAME\_\allowbreak{}OK} y el \texttt{nickname} nuevo.
      \item El SJM actualiza el nombre en la interfaz, muestra la \texttt{GUI\_\allowbreak{}MessageSuccess} y vuelve a la \texttt{GUI\_\allowbreak{}Profile}, donde la opción de cambiar el nombre ya aparece deshabilitada.
      \item Fin del caso de uso.
    \end{cuenum}} \\ \hline
  \textbf{Flujos alternos} & \cuCelda{\textbf{\mbox{FA-01}: El Jugador ya gastó su cambio}\par
    \begin{cuenum}[start=1]
      \item El SJM no habilita el campo y muestra, junto al \texttt{nickname}, la fecha en que se cambió, leída del \textit{HistorialNickname}.
      \item Fin del caso de uso.
    \end{cuenum}} \\
   & \cuCelda{\textbf{\mbox{FA-02}: El nickname no cumple el formato}\par
    \begin{cuenum}[start=1]
      \item El SJM resalta el campo y escribe junto a él qué se espera: entre tres y treinta caracteres, sin espacios.
      \item El flujo continúa en el paso 2 del FN.
    \end{cuenum}} \\
   & \cuCelda{\textbf{\mbox{FA-03}: El nickname nuevo es igual al actual}\par
    \begin{cuenum}[start=1]
      \item El SJM escribe junto al campo que debe ser distinto, sin abrir ninguna ventana: gastar el cambio en el mismo nombre sería el peor resultado posible.
      \item El flujo continúa en el paso 2 del FN.
    \end{cuenum}} \\
   & \cuCelda{\textbf{\mbox{FA-04}: El Jugador cancela en la confirmación}\par
    \begin{cuenum}[start=1]
      \item El Jugador selecciona ``Cancelar''.
      \item El SJM cierra la \texttt{GUI\_\allowbreak{}MessageConfirm} y vuelve al paso 2 del FN con el \texttt{nickname} capturado intacto. No se gastó nada.
    \end{cuenum}} \\
   & \cuCelda{\textbf{\mbox{FA-05}: La sesión ya no es válida}\par
    \begin{cuenum}[start=1]
      \item El Servidor de partidas responde con \texttt{SESSION\_\allowbreak{}INVALID}.
      \item El SJM muestra la \texttt{GUI\_\allowbreak{}MessageWarning} indicando que la sesión terminó y despliega la \texttt{GUI\_\allowbreak{}Login}.
      \item Fin del caso de uso.
    \end{cuenum}} \\
   & \cuCelda{\textbf{\mbox{FA-06}: El servidor detecta que el cambio ya se gastó}\par
    \begin{cuenum}[start=1]
      \item El Servidor de partidas responde con \texttt{CHANGE\_\allowbreak{}NICKNAME\_\allowbreak{}FAILED} y el motivo \texttt{CAMBIO\_\allowbreak{}YA\_\allowbreak{}UTILIZADO}, con la fecha del cambio anterior.
      \item El Servidor de partidas registra el intento en la bitácora del servidor: la interfaz no debería haberlo permitido, así que es indicio de cliente modificado.
      \item El SJM muestra la \texttt{GUI\_\allowbreak{}MessageWarning} con la fecha y deshabilita la opción.
      \item Fin del caso de uso.
    \end{cuenum}} \\
   & \cuCelda{\textbf{\mbox{FA-07}: El servidor rechaza un nickname que el cliente había aceptado}\par
    \begin{cuenum}[start=1]
      \item El Servidor de partidas registra la discrepancia en la bitácora del servidor con la \texttt{version\_\allowbreak{}cliente}.
      \item El Servidor de partidas responde con \texttt{CHANGE\_\allowbreak{}NICKNAME\_\allowbreak{}FAILED} y el motivo \texttt{DATOS\_\allowbreak{}INVALIDOS}, sin gastar el cambio.
      \item El flujo continúa en el paso 2 del FN.
    \end{cuenum}} \\
   & \cuCelda{\textbf{\mbox{FA-08}: El nickname contiene lenguaje inapropiado}\par
    \begin{cuenum}[start=1]
      \item El Servidor de partidas responde con \texttt{CHANGE\_\allowbreak{}NICKNAME\_\allowbreak{}FAILED} y el motivo \texttt{NICKNAME\_\allowbreak{}NO\_\allowbreak{}PERMITIDO}, sin indicar qué término activó el filtro y sin gastar el cambio.
      \item El Servidor de partidas registra el rechazo en la bitácora del servidor.
      \item El SJM muestra la \texttt{GUI\_\allowbreak{}MessageWarning} indicando que elija otro, y limpia el campo.
      \item El flujo continúa en el paso 2 del FN.
    \end{cuenum}} \\
   & \cuCelda{\textbf{\mbox{FA-09}: El nickname ya está en uso}\par
    \begin{cuenum}[start=1]
      \item El Servidor de partidas arma hasta tres sugerencias disponibles a partir del capturado.
      \item El Servidor de partidas responde con \texttt{CHANGE\_\allowbreak{}NICKNAME\_\allowbreak{}FAILED}, el motivo \texttt{NICKNAME\_\allowbreak{}EN\_\allowbreak{}USO} y las sugerencias, sin gastar el cambio (\mbox{RN-05}).
      \item El SJM las muestra como opciones seleccionables.
      \item El flujo continúa en el paso 2 del FN.
    \end{cuenum}} \\ \hline
  \textbf{Excepciones} & \cuCelda{\textbf{\mbox{EX-01}: Fallo al escribir en la base de datos}\par
    \begin{cuenum}[start=1]
      \item El Servidor de partidas revierte la transacción, de modo que no quede el \texttt{nickname} cambiado sin su \textit{HistorialNickname}, que es lo único que impide cambiarlo otra vez (\mbox{RN-04}).
      \item El Servidor de partidas registra la excepción con un identificador de incidencia y responde con \texttt{SERVER\_\allowbreak{}ERROR}.
      \item El SJM muestra la \texttt{GUI\_\allowbreak{}MessageError} con el identificador. El cambio no se gastó.
      \item El flujo continúa en el paso 2 del FN.
    \end{cuenum}} \\
   & \cuCelda{\textbf{\mbox{EX-02}: Fallo al confirmar la transacción}\par
    \begin{cuenum}[start=1]
      \item El Servidor de partidas trata la transacción como fallida, registra la excepción señalando que el estado es indeterminado y responde con \texttt{SERVER\_\allowbreak{}ERROR}.
      \item El SJM muestra la \texttt{GUI\_\allowbreak{}MessageError} indicando que recargue el perfil antes de volver a intentarlo: si el cambio sí se escribió, el campo aparecerá ya deshabilitado.
      \item Fin del caso de uso.
    \end{cuenum}} \\
   & \cuCelda{\textbf{\mbox{EX-03}: Se agota el tiempo de espera de la respuesta}\par
    \begin{cuenum}[start=1]
      \item El SJM cierra la conexión y descarta el intercambio, sin suponer que el servidor no lo procesó.
      \item El SJM muestra la \texttt{GUI\_\allowbreak{}MessageError} indicando que recargue el perfil para comprobar si el cambio se aplicó, y advirtiendo expresamente que no lo intente de nuevo sin comprobarlo.
      \item Fin del caso de uso.
    \end{cuenum}} \\
   & \cuCelda{\textbf{\mbox{EX-04}: El mensaje recibido está malformado}\par
    \begin{cuenum}[start=1]
      \item El SJM descarta el mensaje, cierra la conexión y registra en la bitácora local el contenido truncado.
      \item El SJM muestra la \texttt{GUI\_\allowbreak{}MessageError} indicando que la respuesta no pudo interpretarse.
      \item Fin del caso de uso.
    \end{cuenum}} \\ \hline
  \textbf{Postcondiciones} & \cuCelda{\begin{cureglas}
      \item \textbf{\mbox{POST-01}} Si el caso termina por el FN, el \textit{Usuario} tiene el \texttt{nickname} nuevo y ningún otro \textit{Usuario} puede tomarlo.
      \item \textbf{\mbox{POST-02}} Si el caso termina por el FN, existe un \textit{HistorialNickname} con el nombre anterior, el nuevo y la fecha, y el Jugador ya no puede volver a cambiarlo.
      \item \textbf{\mbox{POST-03}} Si el caso termina por el FN, el \texttt{nickname} anterior queda libre y otro \textit{Usuario} puede tomarlo (\mbox{RN-09}).
      \item \textbf{\mbox{POST-04}} Si el caso termina por cualquier flujo alterno o excepción, el \texttt{nickname} no cambió y el cambio sigue disponible.
      \item \textbf{\mbox{POST-05}} Las partidas, los mensajes y los reportes anteriores siguen apuntando al mismo \texttt{id\_\allowbreak{}usuario}, así que el historial no se rompe pese al cambio de nombre.
    \end{cureglas}} \\ \hline
  \textbf{Reglas de negocio} & \cuCelda{\begin{cureglas}
      \item \textbf{\mbox{RN-01}} El \texttt{nickname} puede cambiarse \textbf{una sola vez} por cuenta. No hay excepción ni forma de reponer el cambio desde la aplicación.
      \item \textbf{\mbox{RN-02}} El \texttt{nickname} mide entre tres y treinta caracteres, es único y se compara sin distinguir mayúsculas ni acentos, de modo que dos jugadores no puedan elegir nombres que se vean iguales.
      \item \textbf{\mbox{RN-03}} Todo cambio queda registrado en \textit{HistorialNickname} con el nombre anterior. Sin ese rastro, un jugador reportado podría volverse irreconocible cambiándose el nombre.
      \item \textbf{\mbox{RN-04}} La escritura del historial y la actualización del \texttt{nickname} ocurren en una sola transacción. Que se aplique el cambio sin el historial daría un cambio gratis; el historial sin el cambio, un cambio perdido.
      \item \textbf{\mbox{RN-05}} Un intento fallido no gasta el cambio. Solo lo gasta un cambio efectivamente aplicado.
      \item \textbf{\mbox{RN-06}} Las cuentas de invitado no pueden usar este caso: su \texttt{nickname} es provisional y el definitivo se elige en el \mbox{CU-07}, que no consume el cambio.
      \item \textbf{\mbox{RN-07}} El cambio exige una confirmación explícita que muestre el nombre nuevo tal cual quedará. Es la confirmación más fuerte del sistema junto con la del borrado de cuenta, porque es igual de irreversible.
    \end{cureglas}} \\
   & \cuCelda{\begin{cureglas}
      \item \textbf{\mbox{RN-08}} Todas las validaciones del cliente se repiten en el servidor.
      \item \textbf{\mbox{RN-09}} El \texttt{nickname} anterior no se reserva. Mantenerlo bloqueado protegería a quien se cambió el nombre de que otro lo tomara, pero a costa de ir agotando los nombres disponibles.
    \end{cureglas}} \\ \hline
  \textbf{Observación} & \cuCelda{\textit{Sobre por qué el nombre se guarda en el Usuario y el historial aparte.}\par
    La alternativa era no tener \texttt{nickname} en \textit{Usuario} y leer siempre el último \textit{HistorialNickname}. Se descartó porque el nombre se consulta en casi todas las pantallas —ranking, amigos, chat, partida— y obligar a una consulta adicional en cada una encarece lo frecuente para abaratar lo excepcional. El \texttt{nickname} vive en \textit{Usuario}, que es donde se lee, y \textit{HistorialNickname} guarda solo lo que ya no está.} \\ \hline
  \textbf{Datos que toca} & \cuCelda{\begin{cudatos}
      \item \textit{Sesion}: lee por token \textit{(paso 8)}
      \item \textit{HistorialNickname}: lee, para comprobar que no se haya gastado \textit{(paso 9)}
      \item \textit{HistorialNickname}: crea \textit{(paso 14)}
      \item \textit{Usuario}: lee \texttt{nickname} de otros, para comprobar unicidad \textit{(paso 12)}
      \item \textit{Usuario}: escribe \texttt{nickname} \textit{(paso 15)}
      \item \textit{PalabraProhibida}: lee el catálogo con ámbito de nickname \textit{(paso 11)}
    \end{cudatos}} \\ \hline
\end{fichacu}
\clearpage
